% preamble-single.tex

\usepackage{float}
\usepackage{pdflscape}
\usepackage{multicol}
\usepackage{xcolor}
\usepackage{svg}
\usepackage{graphicx}

% --- Fallback pour Pandoc ---
% Évite l'erreur \pandocbounded undefined si toutes les images sont dans
% des tableaux (le filtre Lua convertit le tableau en RawBlock LaTeX, 
% cachant les images à Pandoc qui "oublie" d'insérer cette macro).
\makeatletter
\providecommand{\pandocbounded}[1]{%
  {\def\maxwidth{\ifdim\Gin@nat@width>\linewidth\linewidth\else\Gin@nat@width\fi}%
   \setkeys{Gin}{width=\maxwidth,keepaspectratio}%
   #1}%
}
\makeatother

% --- Résumé (page de garde) ---------------------------------------------
% Adaptation pour éviter le conflit avec la classe scrbook[cite: 2].
\let\abstract\relax
\let\endabstract\relax
\providecommand{\abstractname}{Résumé}
\newenvironment{abstract}{%
  \small
  \begin{center}\bfseries\abstractname\end{center}
  \quotation
}{\endquotation}

% --- Définitions de couleurs custom ---
\definecolor{gris}{RGB}{128,128,128}
\definecolor{whitesmoke}{rgb}{0.96, 0.96, 0.96}

% Couleurs CSS de secours pour les spans de surlignement
% ( [texte]{style="background-color: ..."} du Markdown, cf. filtre.lua qui
% met le nom en minuscules avant de l'utiliser en LaTeX ). yellow/olive/
% orange existent deja nativement dans xcolor (ses 19 couleurs de base) ;
% les teintes CSS supplementaires utilisees dans les documents doivent
% etre declarees ici.
\providecolor{yellow}{rgb}{1,1,0}
\providecolor{olive}{RGB}{128,128,0}
\providecolor{skyblue}{RGB}{135,206,235}
\providecolor{lightgreen}{RGB}{144,238,144}
\providecolor{lightsalmon}{RGB}{255,160,122}
\providecolor{lightpink}{RGB}{255,182,193}
\providecolor{royalblue}{RGB}{65,105,225}
\providecolor{forestgreen}{RGB}{34,139,34}

% --- En-têtes et pieds de page (fancyhdr) ---
\usepackage{fancyhdr}
\pagestyle{fancy}
\addtocontents{toc}{\protect\hypertarget{toc}{}}
\fancyhf{}
\fancyhead[LE,RO]{\small\hyperlink{toc}{\textcolor{gris}{|Sommaire|}}}
\fancyhead[LO]{\small\textcolor{gris}{\leftmark}}
\fancyhead[RE]{\small\textcolor{gris}{\rightmark}}
\fancyfoot[LE,RO]{\small\textcolor{gris}{\thepage}}

% --- Styles des titres de niveau 4 et 5 ---
\RedeclareSectionCommand[beforeskip=-10pt plus -2pt minus -1pt,afterskip=1sp plus -1sp minus 1sp,font=\normalfont\itshape]{paragraph}
\RedeclareSectionCommand[beforeskip=-10pt plus -2pt minus -1pt,afterskip=1sp plus -1sp minus 1sp,font=\normalfont\scshape,indent=0pt]{subparagraph}

% --- Redéfinition de longtable pour les grands tableaux ---
\usepackage{longtable}
\usepackage{booktabs}
\usepackage{array}     
\usepackage{multirow}  
\usepackage{calc}      
\usepackage{colortbl}  
\definecolor{grisEntete}{RGB}{217,217,217}
\definecolor{grisBordure}{RGB}{169,169,169}

\let\oldlongtable\longtable
\let\endoldlongtable\endlongtable
\renewenvironment{longtable}{\sffamily\footnotesize\renewcommand{\arraystretch}{1.6}\arrayrulecolor{grisBordure}\addtolength{\leftmargin}{2cm}\addtolength{\rightmargin}{2cm}\oldlongtable}{\endoldlongtable}

% --- Blocs de code et citations (tcolorbox) -----------------------------
\usepackage{tcolorbox}
\tcbuselibrary{breakable,skins}
\definecolor{fondCode}{RGB}{248,248,248}

% Blocs de code avec cadre fin gris et coins arrondis[cite: 2].
\ifcsname Shaded\endcsname\else
  \newenvironment{Shaded}{}{}
\fi
\renewenvironment{Shaded}
  {\begin{tcolorbox}[breakable,enhanced,colback=fondCode,colframe=grisBordure,
     boxrule=0.5pt,arc=2.5mm,left=2mm,right=2mm,top=1mm,bottom=1mm]}
  {\end{tcolorbox}}

% Citations stylisées avec un trait vertical[cite: 2].
\let\oldquote\quote
\let\endoldquote\endquote
\renewenvironment{quote}
  {\oldquote\begin{tcolorbox}[breakable,blanker,
     borderline west={1pt}{0pt}{grisBordure},left=3mm]}
  {\end{tcolorbox}\endoldquote}

% --- Code inline (`...`) : pastille grise --------------------------------
% conf/filtre.lua enveloppe chaque code inline dans \inlinecode. Le
% \fboxsep est reduit dans un groupe : la pastille colle au texte sans
% modifier les autres boites du document. Meme gris que le fond du code
% inline de la Liseuse (--background-contrast-grey).
\definecolor{fondCodeInline}{RGB}{238,238,238}
\newcommand{\inlinecode}[1]{{\setlength{\fboxsep}{1.5pt}\colorbox{fondCodeInline}{#1}}}

% --- Encadres « callout » ( ::: {.callout-note} ... ::: ) -----------------
% conf/filtre.lua emet, pour chaque bloc callout :
%   \begin{calloutbox}{<options tcolorbox>}{<titre>} ... \end{calloutbox}
% Les options portent la couleur du type et, si le Markdown fixe un
% style="width: NN%", la largeur (+ « center » si margin: auto).
% Couleurs alignees sur les jetons DSFR de la Liseuse HTML.
\definecolor{calloutNote}{HTML}{0063CB}      % info      (bleu)
\definecolor{calloutTip}{HTML}{18753C}       % success   (vert)
\definecolor{calloutWarning}{HTML}{B34000}   % warning   (orange)
\definecolor{calloutCaution}{HTML}{716043}   % tournesol (ocre)
\definecolor{calloutImportant}{HTML}{CE0500} % error     (rouge)
\definecolor{lightyellow}{HTML}{FFFFE0}

% Le titre est mis entre accolades : sans elles, une virgule dans le
% titre serait lue comme un separateur de cles par tcolorbox.
\newenvironment{calloutbox}[2]{\begin{tcolorbox}[#1,title={#2}]}{\end{tcolorbox}}

% --- Surlignage de texte ---
\let\ul\relax % Évite le conflit avec le package 'ulem' qui définit aussi \ul
\let\st\relax % Évite le conflit sur \st (strikethrough)
\let\hl\relax % Évite le conflit sur \hl (highlight)
\usepackage{soulutf8}
\sethlcolor{yellow}

% --- Listes et niveaux d'imbrication ---
\usepackage{enumitem}
\setlistdepth{5}
\setlist[itemize]{label=\textendash}  % puce = tiret pour tous les niveaux[cite: 2]

% --- Numérotation limitée ---
% Seuls les titres de niveau 1 (chapitres, "#") sont numérotés.
\setcounter{secnumdepth}{1}

% --- Sommaire (table des matières) ---
% Le sommaire est désormais généré par le mécanisme natif de Pandoc
% (toc: true dans defaults-single.yaml), qui place \tableofcontents
% juste après \maketitle (et le résumé éventuel), donc sous le titre.
% toc-depth: 6 dans le YAML pilote la profondeur du sommaire. On ajoute
% simplement un saut de page après le sommaire pour que le corps du
% document démarre sur une nouvelle page.
\let\oldtableofcontents\tableofcontents
\renewcommand{\tableofcontents}{\oldtableofcontents\clearpage}

% --- Options KOMA diverses ---
% \setkomafont{subtitle}{\normalfont\itshape\Large}
% \KOMAoptions{titlepage=false}